\documentclass[aspectratio=169]{beamer}
\usetheme{Madrid}
\usecolortheme{whale}

\usepackage{amsmath,amssymb}
\usepackage{graphicx}
\usepackage{booktabs}
\usepackage{siunitx}
\usepackage{pgfpages}

% Change to `show notes` when preparing a presenter-notes copy.
\setbeameroption{hide notes}

\newcommand{\safeincludegraphics}[2][]{%
  \IfFileExists{#2}{\includegraphics[#1]{#2}}{%
    \fbox{\begin{minipage}[c][3.2cm][c]{6.2cm}\centering
      \small Missing file:\par \texttt{#2}
    \end{minipage}}%
  }%
}

\title[Project Short Title]{Technical Presentation Title}
\subtitle{Engineering Analysis, Design, and Results}
\author[Hossein Molhem]{Hossein Molhem}
\institute[]{Institution or Company\\Department or Business Unit}
\date{\today}

\begin{document}

\begin{frame}
  \titlepage
\end{frame}
\note{Introduce the problem, audience, and purpose of the presentation.}

\begin{frame}{Roadmap}
  \begin{enumerate}
    \item Problem and motivation
    \item Methodology
    \item Results
    \item Engineering interpretation
    \item Conclusions and next steps
  \end{enumerate}
\end{frame}
\note{Briefly explain the logical flow of the talk.}

\begin{frame}{Engineering Problem}
  \begin{columns}[T]
    \begin{column}{0.56\textwidth}
      \begin{itemize}
        \item Define the physical or computational problem.
        \item State measurable requirements.
        \item Identify constraints and assumptions.
      \end{itemize}
      \[
        y=\mathbf{H}x+n
      \]
    \end{column}
    \begin{column}{0.40\textwidth}
      \begin{block}{Design objective}
        Replace this block with the primary engineering objective.
      \end{block}
    \end{column}
  \end{columns}
\end{frame}
\note{Explain the system model and why the problem matters.}

\begin{frame}{Methodology}
  \begin{columns}[T]
    \begin{column}{0.55\textwidth}
      \safeincludegraphics[width=\linewidth]{figures/system-overview.png}
    \end{column}
    \begin{column}{0.41\textwidth}
      \begin{itemize}
        \item Modeling or experiment
        \item Simulation or implementation
        \item Validation criteria
      \end{itemize}
    \end{column}
  \end{columns}
\end{frame}
\note{Describe the workflow without overwhelming the audience with implementation detail.}

\begin{frame}{Key Results}
  \begin{table}
    \centering
    \begin{tabular}{lcc}
      \toprule
      Metric & Baseline & Final\\
      \midrule
      Error & 8.2\% & 1.4\%\\
      Runtime & \SI{120}{\milli\second} & \SI{24}{\milli\second}\\
      \bottomrule
    \end{tabular}
  \end{table}
  \begin{alertblock}{Engineering interpretation}
    State what changed, why it changed, and why the result matters.
  \end{alertblock}
\end{frame}
\note{Lead with the result, then explain the engineering significance.}

\begin{frame}{Conclusions}
  \begin{itemize}
    \item State the strongest quantitative conclusion.
    \item State the principal limitation.
    \item State the next engineering step.
  \end{itemize}
\end{frame}
\note{Close with evidence, limitations, and a specific next step.}

\end{document}
